\section*{Capítulo \ref{ch:g3:separating-mixtures} --- Separar mezclas}

\begin{solution}{exo:g3:separating-mixtures:1}
El \omterm{def:g3:separating-mixtures:sieve}{tamizado}: los granos de arroz son más pequeños que los guisantes y
caen por un tamiz con los agujeros adecuados. (Separarlos a mano
también funciona, pero despacio.)
\end{solution}

\begin{solution}{exo:g3:separating-mixtures:2}
La tierra baja despacio y forma una capa en el fondo; el agua de encima
se vuelve más transparente. Esto es la \omterm{def:g3:separating-mixtures:decant}{sedimentación}.
\end{solution}

\begin{solution}{exo:g3:separating-mixtures:3}
El \omterm{def:g3:separating-mixtures:filter}{filtrado} es el líquido que atraviesa el filtro. Cuando se hace café,
el \omterm{def:g3:separating-mixtures:filter}{filtrado} es el café de la jarra; los posos se quedan en el papel.
\end{solution}

\begin{solution}{exo:g3:separating-mixtures:4}
Sí. La sal está disuelta, y un filtro no retiene lo que está disuelto:
pasa junto con el agua.
\end{solution}

\begin{solution}{exo:g3:separating-mixtures:5}
\omterm{def:g3:separating-mixtures:evaporate}{Evaporando} el agua: al sol, o con un calor suave, el agua pasa al aire
y el azúcar se queda en el recipiente.
\end{solution}

\begin{solution}{exo:g3:separating-mixtures:6}
La grava se queda en el tamiz. Los granos de arena son más pequeños que
los agujeros, así que caen.
\end{solution}

\begin{solution}{exo:g3:separating-mixtures:7}
(a) \omterm{def:g3:separating-mixtures:sieve}{Tamizado}. (b) \omterm{def:g3:separating-mixtures:filter}{Filtración} (o \omterm{def:g3:separating-mixtures:decant}{sedimentación} y \omterm{def:g3:separating-mixtures:decant}{decantación}).
(c) \omterm{def:g3:separating-mixtures:evaporate}{Evaporación}.
\end{solution}

\begin{solution}{exo:g3:separating-mixtures:8}
Cuatro pasos: rejas, depósitos de \omterm{def:g3:separating-mixtures:decant}{sedimentación}, filtros de arena y
eliminación de microbios. Los depósitos de \omterm{def:g3:separating-mixtures:decant}{sedimentación} eliminan el
barro.
\end{solution}

\begin{solution}{exo:g3:separating-mixtures:9}
Filtrar; recoger la arena del papel; \omterm{def:g3:separating-mixtures:evaporate}{evaporar} el \omterm{def:g3:separating-mixtures:filter}{filtrado} para obtener
la sal.
\end{solution}

\begin{solution}{exo:g3:separating-mixtures:10}
No. La sal del agua del mar está disuelta, y un filtro no retiene lo
que está disuelto: el \omterm{def:g3:separating-mixtures:filter}{filtrado} sigue siendo agua salada.
\end{solution}

\begin{solution}{exo:g3:separating-mixtures:11}
Verter el té con cuidado una vez que las hojas se han depositado en el
fondo (\omterm{def:g3:separating-mixtures:decant}{decantación}), o verterlo a través de un colador o de un filtro
(\omterm{def:g3:separating-mixtures:sieve}{tamizado} o \omterm{def:g3:separating-mixtures:filter}{filtración}).
\end{solution}
